% TOP_PROBLEM: {"id": 30007000, "title": "Finite central dimension of algebraic finitely generated division rings", "category_id": 4, "category": {"id": 4, "name": "algebra", "display_name": "Algebra", "description": "Group theory, ring theory, field theory, and algebraic structures.", "slug": "algebra", "order_index": 4, "created_at": "2026-07-31T15:26:25.671Z"}, "status": "open"}
% FIRST_POSTED: 2026-10-06
% ATTEMPT_STATUS: solved
% OPENAI_CLAIM: {"schema_version":1,"source_repo":"https://github.com/openai/math","source_commit":"adc7f1241b42e322a6451854ab7e4b4c146bf78a","release_date":"2026-10-06","match":"full","resolution":"disproved","independently_reviewed":false,"second_pass":"agreed","families":[{"family":"201","title":"A counterexample to Kurosh’s division-ring problem","subject":"Algebra","manuscripts":[{"title":"A Counterexample to Kurosh’s Division-Ring Problem","pdf_path":"preprints/A-Counterexample-to-Kuroshs-Division-Ring-Problem-September-23-2026/paper.pdf","readme_path":"preprints/A-Counterexample-to-Kuroshs-Division-Ring-Problem-September-23-2026/README.md","date":"2026-09-23","theorem_ref":"preprints/A-Counterexample-to-Kuroshs-Division-Ring-Problem-September-23-2026/build/sections/introduction.tex:19, result labelled thm:main (pinned adc7f1241b42e322a6451854ab7e4b4c146bf78a)"}],"lean":null,"reasoning_trace":null}]}
\documentclass[11pt]{article}
\usepackage[margin=25mm]{geometry}
\usepackage{fontspec}
\setmainfont{Latin Modern Roman}
\usepackage{amsmath,amssymb,mathtools}
\usepackage{unicode-math}
\setmathfont{Latin Modern Math}
\usepackage{enumitem}
\usepackage{xurl,hyperref,bookmark}
\hypersetup{colorlinks=true,urlcolor=blue}
\setcounter{secnumdepth}{0}
\setlength{\parindent}{0pt}
\setlength{\parskip}{5pt}
\setlength{\emergencystretch}{3em}
\begin{document}
\section*{Finite central dimension of algebraic finitely generated division rings}

\subsection{OpenAI's claim}
OpenAI presents a counterexample in the selected manuscript. An algebraic division ring is claimed to be generated as an algebra by two elements while having infinite dimension over its actual center.

\subsection{Outline of the argument}
A countable schedule imposes algebraic relations on successive polynomial expressions. Formal-series lifting inside finite-dimensional division-algebra models retains two independent deformation directions; trace constraints and conjugation calculations preserve generation, and a separate descent argument prevents the final center from making the dimension finite.

\subsection{Scope relative to this problem}
The theorem gives D=F[x,y], every element algebraic over the center F, and [D:F]=infinity. Its two algebra generators meet the finite generation hypothesis and refute the required finite-dimensional conclusion. Within these assumptions, this is a claimed negative answer to the stated question.

\subsection{Formal verification}
Family 201 has no Lean scope page in this release. Its selected result is not presented as formalized. This site has not executed the Lean code. The scope descriptions do not certify the remaining manuscript arguments.

\subsection{Status caveat}
OpenAI attributes these manuscripts to work by an internal model. Its release cautions that arguments without formalization can contain mistakes. This catalogue has not independently checked the proofs; the labels report the mathematical claims attributed to OpenAI.

\subsection{Sources}
\begin{itemize}
\item OpenAI, \emph{A Counterexample to Kurosh’s Division-Ring Problem}, 2026-09-23; citation key \texttt{OAI:A-Counterexample-to-Kuroshs-Division-Ring-Problem-September-23-2026}. \href{https://github.com/openai/math/blob/adc7f1241b42e322a6451854ab7e4b4c146bf78a/preprints/A-Counterexample-to-Kuroshs-Division-Ring-Problem-September-23-2026/paper.pdf}{PDF (GitHub)}; \href{https://github.com/openai/math/raw/adc7f1241b42e322a6451854ab7e4b4c146bf78a/preprints/A-Counterexample-to-Kuroshs-Division-Ring-Problem-September-23-2026/paper.pdf}{PDF (direct)}; \href{https://github.com/openai/math/blob/adc7f1241b42e322a6451854ab7e4b4c146bf78a/preprints/A-Counterexample-to-Kuroshs-Division-Ring-Problem-September-23-2026/README.md}{Manuscript page and citation}; \href{https://github.com/openai/math/blob/main/preprints/A-Counterexample-to-Kuroshs-Division-Ring-Problem-September-23-2026/paper.pdf}{Latest version}.
\item \href{https://github.com/openai/math/blob/adc7f1241b42e322a6451854ab7e4b4c146bf78a/CONTENTS.md}{OpenAI catalogue entry}.
\end{itemize}
\end{document}
